% Extended lesson notes for the middle group of problems.

\expandafter\def\csname sourceexcerpt@26\endcsname{%
Kadane's algorithm we really are just scanning throughout
the rest of the subarray to find the maximum sum accordingly depending on various
combinations and it is a pretty efficient algorithm to be able to solve such a problem
at the same time and so what I will be doing in this case is use a
new type of algorithm called Kadane's algorithm And in this one what we will be doing
is that we are going to have a series of steps that will help us find our particular
subarray.

For the input that we have is a given array and what we need to do is return the
maximum sum of a subarray. What are we going to do with the data? We are going to find
we're going to scan through the array to find the maximum sum of a subarray. We should
be able to create a subarray while we're scanning through while we are scanning it
through our given array. To start with, we're going to use Kadane's algorithm and we
will we're going to scan through the array and create a subarray to find the maximum
sum of the subarray.

ultimately to find the maximum sum is what we are returning. So to be able to start we
are first going to initialize a current max number and then a we're going to
initialize a global max number which should be an index of the array next of the given
array. should be the initial. So that it's much like in the case of having a variable
and setting it to zero. Okay.

So to be able to start an current max is going to be equal to nums of zero as we are
already given a vector at the top and then we're also going to have a global max as
well of zero and then what we will be doing is that we are one going to be scanning
through array. So we're going to have a for loop int i is going to be.
}

\expandafter\def\csname sourceexcerpt@27\endcsname{%
The reason being is because after scanning every single possible subarray to start
with the x two and three happen to be the highest relative to three and two and then
-2 and four to start when comparing all the rest of the numbers all together and then
we can also have an edge case at the same time as well. So given that we have -2 and1
the result cannot be two because -2 and 1 is not a subarray. So they all have to be
connected one way or another.

So we know the relationship as result from coming from set examples and the technique
that we are going to focus on is not going to be a data structure to start with. it's
going to be focusing on the STL library in C++. So we're going to do a swap and max
technique where these methods are from the STL library where we're able to swap some
of the variables that we're going to initialize first then we are going to loop
through arrays. So what we are going to do is that we will initialize our min, max and
answer relative to the index of the array given array.

And while we are scanning through the array, we have to take in mind of some edge
cases because as highlighted in the example, we will have negative numbers to start
with and the constraints given the fact that numbers of I will range from -10 all the
way to 10. consider negative numbers. That is why we're going to use the swap. But
really our main technique is going to be the max to begin edge case of negative
numbers. And then what we are going to do is also max and minimum calculation and swap for
edge case. Okay.

So overall the technique that we're going to focus on is focusing on re redeclaration
from the STL library. We're going to redeclare our initialized variables first for
both the max and the minimum as they are going through the rest of the array because we
have to consider all other possibilities when we are scanning through the array and we
will declare the min and max products accordingly.
}

\expandafter\def\csname sourceexcerpt@28\endcsname{%
Okay, I'm just going to put in the rest of comments over here to start. Let's get
further out. All right, you're given a string as an integer K. You can choose any
character of the string and change this to other uppercase English character. You can
perform this operation at most k times return the length of the longest substring
container the same letter you can get after performing the above operations. Okay, so
just to give a breakdown of what we're doing. So it looks like that we are will be
returning the length of the longest substring containing the same letter you can get
after performing the above operations.

Now the same amount. So in particular by what they mean to be able to start is that we
just need to return the maximum length of the substring that can be performed by
replacing as many most characters. So going over example one we have a given input of
string a b a ab a b and k of two.

Now and then the output of this is for and the reason being as they state is replace
the two A's or two B's and vice versa as in how many times can you replace the number
of characters inside of our given string and in this case for example a b a bba it's
asking us to replace one of the characters and in this output it said we can replace
it four times such as replace the one a in the middle of b and form a b BB B A with
the longest substring of being BBB with the longest repeating character.

So after we replace the character we would have to see what the longest substring
would be after we replace for instance two characters and then we could have the
longest output accordingly as a result from replacing the two characters. Okay. So we
are so k represents the number of characters we can replace and then we are again
returning the length of long substring taking the same number you can.
}

\expandafter\def\csname sourceexcerpt@29\endcsname{%
This problem uses dynamic programming. Dynamic programming solves a complex problem
by breaking it into smaller overlapping subproblems and saving answers that will be
needed again. Instead of solving the whole string problem at once, we first describe
the answer for a smaller range and then build larger answers from it.

So it's taking the elephant bit by bit and then you just expand upon the relationships
over variables that you have already been given to start with and what variables you
can also keep track of. So what we are going to create we are going to create two what
we are going to do is that we will need to think of the string and we would need to
keep track of what we what we have already considered as well.

So the first thing that's already going to be keep kept in mind is that we will create
a vector where it will be a matrix to keep track of all the characters visited and if
they are matching or not.

not the matrix will be called DP with N * N and then what we are going to be doing is
that for what we have already initialized when expanding upon the relationship we're
also going to have start and max length of our string and then to what we are doing
with the following data just expanding upon it is that we will first loop through the
string s from the end to start because we're going to break up the string into two
parts and then in the next part we will from end to start and in the same operation we
will loop through the string s from the beginning to the end and then we're going to
check for every single and we're going to check for every single character which has
been visited and then visited and we're not necessarily storing since.
}

\expandafter\def\csname sourceexcerpt@30\endcsname{%
The answer is B with the length of one. Hence right over here. Pretty straightforward.
And for this one, P WW K W with the output of three. The answer is W K with a length
of three right here. And notice the answer must be a substring and it cannot be pwke.
Hence, because we have another character that is in between PW and K. So with our
following output we have here an S contains English. Okay. So be able to break down
over the questions we have already answered. So we already know what our input's going
to be from our question.

You can even look at the parameters at the same time to discern what our input's going
to be. It's a string. And then we already know what we're going to do with the data.
We need to find the length of longest substring without repeating characters. And for
our output as highlighted by the example, we will return an integer representing the
number of unique elements we have visited in our string. What is the longest substring
with unique characters? Those are our questions that we would have.

And so the technique we will useday is going to be called the
sliding window technique where unlike binary search where it's divide and conquer
method and you have three different variables of a low, highend, and a middle and then
you find one area and then cut off another area until you found the specific element
you're looking for. And it's not the two-pointer technique where you just have two
variables, a low and a high end where in your sliding window you are given a range
that you also need to figure out within a data structure you're trying to scan through
like an array.

The sliding window technique has one two points and then also has a length that one
could already set. So there are three variables one must construct particular low,
high and maximum length for the sliding window technique. And with our sliding window
technique, we will scan through our array to keep track over the number of unique
characters we have visited through our array and how we're going to keep track of
everything.
}

\expandafter\def\csname sourceexcerpt@31\endcsname{%
So return the number of combinations of palindromic substrings and a palindrome is as
stated before is when it reads the same forward and backwards. It's like damn it I'm
mad spelled backwards as racecar like same way the same amount of characters
in the string accordingly forward and back. All right. And so this type of problem for
palindrome dynamic programming is a popular technique that we would use for our
solution. And I'm just going to go further in depth inside the steps of what we got to
do. So break it down.

First thing we need to do is that we need to initialize the string length and the
matrix that will have the string length on the inside to check to see if there is a
palindrome. if so it will be true. This may be a little bit confusing and for the
first step but what I mean by the first step to start is that we are going to be
creating a matrix and I will write this out after we go through the steps.

So the first step we will do is that I'm going to have n be equal to s dolength right
here and then we're going to create our matrix over here. We're going to initialize
our matrix where it's going to be a vector bull. It's going to be a matrix accordingly
where it is n which is the size of our string vector boolean and false right here. And
this will represent the and n represents the size of our given string that we have.

And we're checking to see by using this matrix throughout each of the steps that we
will be doing because we're going to be looping through our entire string to check for
palenromes. We'll be using this to keep track of if there is a palrrome or not. Okay.
So for the second step, what we need to do is that we need to create a count as well.
And then for the third, we're going to iterate through our string s two.
}

\expandafter\def\csname sourceexcerpt@32\endcsname{%
So interval we have array of intervals as even indicated by right up here in the
parameter. So let's go over the example. Example, we have interval of 0 3 0 30 5 10 15
20 where intervals I equals start I and I. Determine if a person could attend all the
meetings. Well, the question is can the person attend all Okay.

So to be able to give a rundown of this whole thing and to be able to break up the
question is that so if we have intervals in particular so the med is going to be
starting at this time and then ends at this time starts at this time ends at this time
starts at this time ends at this time as a result from this given the fact it's
overlapping the other times therefore it's going to be false.

But then for example two we have seven 10 two and four where it looks like it's just
for three hours over here and looks like it's two hours over here they don't overlap
so it's going to be true. Okay.

So to be able to go further down over what we will use the data structure we will be
using is going to be a min-heap which is built using a priority queue or that is the
principle of last in first out much like in the grocery no it's first in first out
sorry on that priority queue or Q's in general is first in first out much like in the case
of the grocery store for instance and shopping cart. That's the best way of
metaphorizing it. And so what we will be using this to sort the intervals. All right.
Here we go.

And then we need to ask ourselves what to do with the data. Okay. So what we need to
do with the data is first sort the intervals based on their start times. Then we will
iterate over the sorted intervals adding each interval to priority queue and heap. Okay,
just clean this up. All right. If the start time of the current interval is greater
than or equal to the smallest end time, the root of the min-heap, we can use the same
room and pop the end time from min heat.
}

\expandafter\def\csname sourceexcerpt@33\endcsname{%
So going over the example, example one, is it possible to jump from one index jump
from the very first index and all the way to each element accordingly. Okay. And each
element in the array presents your maximum jump position. Okay. right here. Okay.
Yeah, right there. That's pretty straightforward. So, we have this three zero false.
You can't because it's zero. Makes sense. All right. And the jump we're going to for
output is going to be return true if we can make the jump. false if we cannot. So
looks like for this entire thing we will be using the greedy algorithm and what are we
going to do with the data?

Well, three four five. All right. So one we need to Overall, traverse the array from
right to left and maintain a last position which is initially the last index and then
the current index plus the value at the index. X is greater than or equal to the last
position, we update the last position to be the current index. So how are we going to
do this? Well, based on the number of steps we are going to do, which is going to be
four, we're going to initialize the last position to be the last index of nums. And
then we are going to traverse through the array from left to right.

And then if the current index Plus the value at that index is greater than or equal to
the last position. Update last position to be the current index. So that way we know
exactly if we have attained our maximum jump. And then finally if and finally we will
return true if the last position equals zero otherwise false. So we're going to first
do our last position equals nums do size minus one over here. Then we're going to do
step two. we are going to iterate through our entire array. Zero is going to be equal
to nums do size minus one since we're going from the rear equals zero.

and then I minus over here. Then step three if I like for example hang on I'll write
this out and be able to showcase what we have.
}

\expandafter\def\csname sourceexcerpt@34\endcsname{%
Okay so we are given a tree and we need to return maximum depth and so the technique
we're going to be doing is depth-first search hence in the name it's useful to find depth
and then there are typically two techniques in trees breadth-first search and depth-first search
deer search utilizes more of a stack buffer search uses a Q in this case but C++ and
their STL library is very intuitive when it comes to their own library and so we'll
just be focusing on leveraging on that and it's really four easy steps to go over. So,
it's going to be really fast.

So, the first step we're going to do, we're going to have our edge case when we don't
have when don't have any value in the root is equal to null. Second step we will have
is calculate the maximum depth of the left subree. Then we'll do the same for the
right finally return the max of depth of both the left and right sub trees plus one to
include the root. All right. So this will be straightforward.

So if root equals null ptr we return turn return turn return turn return turn return
turn return turn return turn return turn return turn return turn return turn zero
since we don't have anything in here and we are returning a number and then step two
step two the maximum so it's going to be equal to the max depth root left. We'll do
the same thing for the right stuff. Just going to make sure we have everything already
set and it is right. All right, passes all the test cases and here we go. Woohoo. That
was quick and efficient.

Okay, just be able to give a rundown and breakdown of everything we already did. So
the time collective for our following solution is going to be O of N where N is
representing the number of nodes in the tree to start and that we are finding both the
right and left depth of our following tree. We have to keep track of all of the nodes
that have been visited to be able to get our max step.
}

\expandafter\def\csname sourceexcerpt@35\endcsname{%
So we would have we have this accordingly and we are going to sort them out first. So
we have one two three and four and it's asking us the smallest element throughout the
rest of our binary search tree that we need. It's just all right and our output for
this one will be one accordingly because the value that we have at one that's the
smallest within this area is one. Now if we were to go over the second example, if we
were to have root for this following binary search tree and then we have K.

What we need to do for this specific problem is that we need to first sort everything
out. Four, five, six, and then we also have null and null. But really we're just
focusing on the following numbers right here. Those will be gone. So what we need to
do is that we need to find the third smallest element accordingly and so our output
would be three. That is what they mean by the kth smallest value at the following
indexed.

So what we need to do is that we need to find a way or we can one sort all the values
within our binary search tree and then according to K find the k smallest value which
is indexed at the given position. So if we were to have for instance k be four for
example the solution would be four. But our biggest problem that we have is that we
have a binary search tree rather than a sorted array. The useful technique is an
inorder traversal. It begins at the leftmost node, visits the current node, and then
moves through the right subtree. For a binary search tree, this visits the values in
increasing order.
}

\expandafter\def\csname sourceexcerpt@36\endcsname{%
So if we have our input right over here root is equal to 6 280479 and then null 3 and
five then we have P equal to two and then Q equal to 8 lowest common ancestor of nodes
two and six two and 8 is going to be six over here given the fact we are working with
a binary search tree but then if we have the root of 6 28 8 0479 null 35 P and Q 2 and
4 two over here and then we have four accordingly what is the lowest one that we would
have this relation to well according for this one two the long the lowest common
ancestor nodes two and four is two since a node can be a descendant of itself
according to longest to the lowest common ancestor definition.

Okay. So two and one accordingly. All right. So it looks like that the lowest common
ancestor refers to the root at the highest level possible. So meaning in particular
that if we have for example six over here P and Q 2 and 8 they are both descendants of
six. The definition states the lowest common ancestor is defined between two nodes and
the lowest node in T. That means that P and Q are descendants of it. And if four is a
branch of two that means two will be the lowest common ancestor. Okay. So we need to
replicate this behavior for the lowest common ancestor. Just clarifying on this.

And so what we will be doing for our following steps accordingly is that from this we
will be extrapolating on binary search tree properties. on the relationship of the
nodes we have. And in doing so, we will traverse the tree and compare the values of
the current node with the values of P and Q. All right. And so just to be able to go
over the steps, we're going to first have our edge case that we need for the
following.

It looks like for the tree, what we would need to do is that well since we already
given a binary search tree, let me see if there are any other constraints as well.
}

\expandafter\def\csname sourceexcerpt@37\endcsname{%
So are given an array of interval times we need to return number conference rooms
required. And in this case, we're going to use a min-heap, which is a pri which acts
like which is a priority queue to do with the data. So to be able to figure out what we
will be doing with the data, we're going to break it up into four different steps.
One, we are going to sort the intervals by their start times. Second, We will
initialize a min-heap to keep track of all the end times of the meetings. Then three
for each interval in the sorted list of intervals.

If min-heap is not empty and start time of the current interval is greater than or
equal to the smallest end time root of min-heap. Gonna pop the end time from min-heap
and then we will add the current intervals end time to the min-heap. Then we're going
to output put size of the min-heap. That's pretty much it. It's going to be asking
just overall based on the following times how many meetings there are in total. the
minimum number of conference rooms required.

So it could be intervals for instance on if we have two meetings going on at the same
time within a given context like for instance 5 and 10 fall in 0 and or 15 and 20.
We're going to need at least two conference rooms for the star. And in this case, we
can have one to begin with. So, let's kick things off. Step one, go to sort again.
intervals end just calling from the STL library. We're going to create a priority
queue. three. Then we're going to look through part A. If not, min-heap if the
min-heap is not empty.

and the interval at the very first index is greater than or equal to the min-heap at
the top. We're going to min-heap going to pop the heap and then we're going to add the
current element with the min-heap. And then finally we're going to output the size.
So, we're going to turn min heaps size.
}

\expandafter\def\csname sourceexcerpt@38\endcsname{%
So we have all of the numbers inside of our intervals array. So we have 1 2 3 and four
for the most part. And what it looks like, it looks like if we were to remove this
overlapping interval that we have, one and three, because if we were to scan the rest
of the numbers, we see 1 and two is inside of one and three. And two and three is also
inside of 1 and three. So highlight the numbers to visualize the relationship between
all the elements.

And for our input and output and if we were to go forward on example two we have so in
this case it makes sense over exactly why the output would be two in this case for the
minimum number of intervals we have to ultimately get an array that would have
intervals where all the numbers are non-overlapping and we need to return the
number of intervals equals one needs remove to make the rest of the intervals not
overlapping. And for example, three looks like between 1 and two, two and three, zero.

Because even though it may appear at first, if we were to expand more at least on this
example and our output is zero. The reason being is, hey, we have one, two, and three.
One and two, and two and three. Even though they're connected by two, they're not
overlapping. Technically, it's really just in the case like this one for example, one
through three and then we have one and two where we have a number inside of our
interval. Over overlappinging is when we have a number inside another number inside
our interval.

And so the technique that we will be using that we use for
our solution to be able to solve is going to be a greedy algorithm with sorting. And
what I mean by that is that a greedy algorithm is more of a technique really that
explores all the possible relationships within a given series of elements that we are
given. Meaning that if we for instance are given an array, we need to discern exactly
the relationship from all of the elements and all the possibilities that are possible.
}

\expandafter\def\csname sourceexcerpt@39\endcsname{%
And so if we were to break down the example 1 and two, let's say for instance we have
intervals equal 1 3 2 6 8 10 15 and 18. Our expected output that we have is that we
need to merge all the overlapping intervals and return rate of non-overlapping
intervals. So we have 1 6 8 10 15 8. And it looks like from what we have, we're able
to merge the rest of the intervals that we have. And you see 158, 810, and then 1 and
6.

And it may be confusing to be able to start, but we merge 1 and 3 and 2 and 6 together
to make one interval. Since intervals 1 and three and two and six overlap, we merge
them into one and six. Pretty vivid example. right over here. Just going to put this
in the comment section as well. And then we also are able to do the same for intervals
1 and four, four and five. What they come together to become one and five. And it's
because 1 and four and four and five are considered overlapping. All right, it's
pretty self-explanatory or what we need to do.

And so the technique that we will use is we're going to be
doing a sorting approach. with an iterative method where what we will be doing with
the data in order to merge all the inter overlapping intervals is that we will first
sort the intervals based on their start times. Then we will iterate over the sorted
intervals merging the over overlapping intervals and iteration and iterating means to
go through or scan through. That's just a fancy way of putting it.

And so just to be able to go over the steps over how we're going to do this is that
for step one, we need to sort the intervals by their start time. Then step two, we
need to initialize a new list for the result. Then for step three, we're going to do a
for each loop for each interval in the sorted list of intervals.
}

\expandafter\def\csname sourceexcerpt@40\endcsname{%
And so it looks like from the following data over exactly what we need to do. We need
to insert the new interval into the new into the intervals such that intervals is
still sorted in ascending order and they are and all the intervals are nonover
overlapping. All right. And it looks like for our output we need to return another
interval where the numbers are not over overlapping each other with the new inserted
interval we are given. Okay.

And so if we want to break down the following example, say for instance we have
intervals equal 1 3 6 and 9 and we have a new interval 2 and 5 and our output is 1 5 6
and 9. It looks like from the following since the problem also states that we may need
to merge overlapping intervals if necessary because we need to return the new
intervals that don't have any overlapping intervals. So we have an intervals array and
then for our output 1 5 6 and 9 these numbers are non-overlapping.

So the relationship that we are seeing between the input and the output is that well
if we are given a new interval we need to insert that new interval on the inside. So
we would have 1 3 6 and 9. But since 2 and 5 are overlapping 1 and 3. This is why we
merge 1 and 3 with 2 and 5. So we get 1 and five accordingly and then 6 and 9 because
all these intervals are non-overlapping and we insert the new interval. Now if we
were to look over example two, we have an array of intervals to start with and a new
interval 4 and 8.

Well, it looks like four and eight merge within three, five, six, seven, and eight and
10. So, these three different intervals are having four and eight be overlapped inside
of them and we need to merge them accordingly, which is why we get three and 10. And
so, our output would be 1 and 2, 3 and 10, 12 and 16 as non-overlapping
intervals.
}

\expandafter\def\csname sourceexcerpt@41\endcsname{%
Okay. So if we want to highlight this, we got ones here, ones there, and they're all
adjacent ones and ones accordingly. We have one island and our output is going to be
one. Now, if we were to expand on example two, we have the grid already here and then
we have ones. So, we got one island, then we have one island over here, then we have
another island over here. The rest of them are just zeros and the output is going to
be three.

So the technique for us to be able to use to accomplish this pattern, we're going to
use depth first search where unlike breath first search, depth-first search and I'll also
have a visual to be able to add and explain more. Deer search focuses on the levels in
particular starting from the leftmost all the way to the right. So it's going down
then it's moving up from each branch accordingly until all the nodes are accounted for
and they're all connected starting from the root. That is the key difference between
breadth-first search and breath research accordingly.

And so the steps needed to be able to solve the following problem is that to keep in
mind over what we're going to do with the data, we're going to we're going to overall
we will traverse the grid and perform depth-first search for each one cell and increment our
counter accordingly. We will create a recursive function called DFS which we will use
to implement depth-first search behavior.

And so just to give a breakdown on steps on how we're going to do this accordingly is
that the first thing that we need to do is that we will steps we're going to
initialize a counter to count the number of islands and that will be int counter zero.
int counter is equal to zero. And then what we will be doing is that we're going to
loop over each cell in the grid every row and column which is going to be indication.
We're going to have a nested for loop.
}

\expandafter\def\csname sourceexcerpt@42\endcsname{%
So looks like from our input it will be a number of nodes in the graph which is going
to be n and a list of edges representing the connections between nodes. And it looks
like our output is going to be the count of connected components in the graph. And it
looks like since we are focusing on connectivity in this case and we are dealing with
a graph problem.

a thing to note in trees popular methods to be able to solve the following solution
would be deaf research or breath research and deaf research would be ideal for a
solution since we are focusing on connectivity of the nodes and the differences
between deaf research and breath research is that deaf research focuses on the
connectivity starting at the root of a given tree or graph and then goes to the
leftmost type of node on the downward word, hence depth in the name. And then does the
same thing for the right branches accordingly starting from the left most branch.

While breath for search on the other hand focuses on the nodes per level. So for
instance on breath for search we have one level and then it just counts the nodes from
the left to the right accordingly. Unlike depress search and since we're focusing
connectivity depth for search would be the technique to use count the number of
connected components. So we will use depth-first search to find the connectivity of the
nodes and count the number of connected components. And what we will do with the data
is that by principle we will create an adjacency list to represent the graph.

Then starting from each unvisited node, we will perform depth-first search to explore all
connected nodes and mark them as visited. And so the steps for us to be able to go
over this entire thing, and I'll also have expanded pseudo code as well, is that for
step one, we're going to initialize an adjacency list to represent the graph. In step
two, we will create a variable to store the count of connected components.
}

\expandafter\def\csname sourceexcerpt@43\endcsname{%
Okay, so to be able to give a breakdown over exactly what we need to do for the
following. the first thing that comes to mind is that well for one our given input is
that we are given a list of nums. especially when looking at the example. And the
expected output is to return the length of the longest consecutive element sequence.
So if we were investigating this, the longest consecutive elements inside of the whole
thing depending on how it's jumbled up at the same time. So we have 1 2 3 and four are
the longest consecutive. So therefore it would be four.

And then since we have all these numbers already jumbled up together, but we still
have an output of nine, we have the following longest consecutive sequence being nine
since it starts from zero and goes all the way up to eight accordingly. And the
technique that we will be using for this is that we will be hashing. Now hashing is a
way where we are able to keep track of the length of a given sequence accordingly
where we are able to break up our given data and keep track of each part accordingly.

So what we will do namely is that we will insert all numbers into a hash set for fast
lookup and then iterate through each number and check if it is the first number of a
sequence. then count the longest sequence it can form. All right. Now, for us to be
able to go over the steps on how to replicate this type of hashing behavior is that
the first step we need to do is that one we need to create a hashet and add all
numbers to it. Then we need to initialize a variable to track the max length of a
sequence.

Then we will loop through each number in the array since we were dealing with arrays.
And then we're going to check if it is the first number of a sequence. num minus one
does int doesn't exist in the hash set and then if it is then count the length of the
sequence it can form and update the max length if necessary.
}

\expandafter\def\csname sourceexcerpt@44\endcsname{%
Okay. Well, it looks like we got to merge the two list and one sorted list. And we
found a way just from the question to be able to break it down based on keywords and
how and knowing how to organize it. So it looks like in particular that we are just
merging all the link lists together. Like we have list one, list two, all of them
being sorted and merged accordingly. If we have an empty list, we just have an empty
list return. And then if we have an empty list and zero. So we have zero right here.

And the numbers can go from -100 to 100 in value for each node. And the number of
nodes in both lists is in the range of 0 and 50 the value. Interesting. And then we
have both list one and list two are sorted in nondereasing order. All right. So, we
already have that. So, luckily what we will do because of how easy it was to be able
to dissect the problem is that we're going to use merge algorithm of merge sort where
as the name explains what we're doing is we will be creating an algorithm to merge
both of the link lists accordingly and then we will be sorting them.

ultimately we can sort them first and merge them all together. Well, since we already
have two sorted link lists to begin with and what we will do is that we're just going
to check while we are scanning through them. So the steps for us to be able to operate
on this are the following. Step one, we need to initialize a new head and current
pointer. Then step two, while both list one and list two are not null, if list one's
value is less than list 2's value, we will add list one's node to the current list and
move list one forward.

Otherwise, we will add list 2's node to the current list and move list two forward.
Then, while we are doing this inside the while loop, we're going to move the current
pointer forward after adding the node since everything is already sorted and that
helps us a lot.
}

\expandafter\def\csname sourceexcerpt@45\endcsname{%
So that means in particular what we will do is that we will for one we would need to
create a structure to compare the values of each node from one list to another. That's
just a key indication over what we need to do for Q's. In our case, what we will be
doing for this data structure, cues in general operate on the principle of a last in
first out situation. Much like in the case of a grocery cart line where people are
going in throughout the line and then they happen to appear out of the line that those
how in some retrospect that's how cues operate.

And so the steps that we would need in order for us to be able to go through are the
following. one, we're going to create a boolean function called compare to compare the
values of each linked list like list one and list two. Second, we're going to
initialize a priority queue that orders the nodes by their values.

the third step and then for each list in the input if the list is not empty we're we
will add our first node to the que then the fourth we will initialize a new head and
current pointer for us to return then we will scan the que and while the Q is not
empty. We are going to move the current pointer forward after adding the node and then
we will return the new list from the head which will be called dummy. So the way how
this is going to work out accordingly are the following steps. Step one, step two,
step three, step four, step five, step six, step seven.

A thing to note since we're going to be comparing the values, we're going to create a
separate function all the way at the top and we're going to start comparing each of
the values accordingly. So the way how it will look is that this is how we create a
new function that will be called later. We have a structure called compare. And what
it's going to be is that it's going to be a boolean function or really just an
operator that will incorporate the parameters of a list node.
}

\expandafter\def\csname sourceexcerpt@46\endcsname{%
And we need to return the length of longest common subsequence. So that's the first
one. return the length of their longest common subsequence. And the technique that
we'll be using for this following solution, if we're going over the examples, namely
what we're trying to do is that we're trying to find the longest common subsequence
between the two patterns of two different texts when we're scanning both of them
accordingly. And the technique we're going to use is a dynamic programming to scan
both text one and two.

Using a dynamic programming table to check for character and a size of our text
strings and the steps for us to be able to go about the this are for the following. So
step one, we're going to initialize a DP table with the size of text one dot size + 1
* text 2 dot size + one and when we do so this is how we create our table. Then what
we will do is that for step two, we're going to loop through each character in both
strings.

And then for step three, if the current characters in both strings are the same, the
value of the current cell is 1 plus the value of the cell diagonally of the table both
above and to the left of it. And for step four, if otherwise if they are different,
the value of the current cell is the maximum between the cell to the left and the cell
above in our dynamic programming table. Finally, for step five, we're going to return
the value in the bottom right corner of the dynamic programming table we created to
show the longest common subsequence.

And I'll also have a visual as well for the dynamic programming table while we are
constructing this whole behavior. now just to clean up everything else and just go
straight into the code over how it will look like. So I'm just going to write the
steps just five steps accordingly. Step one, step two, step three, step four, step
five. Let's just get after this.
}

\expandafter\def\csname sourceexcerpt@47\endcsname{%
Since the largest window of s only has one a and we return an empty string. Okay, so
those are some things to be able to keep in mind. Now, my use since we were dealing
with two strings in particular, the technique that we'll be doing is going to be the
sliding window technique. I'll have a visual on the side to be able to help illustrate
the idea that we're trying to construct for this solution. But the sliding window
technique, we will create two maps to keep track of the frequencies of our characters,
including duplicates.

And maps the best way of visualizing them is that they have a key value pair type of
relationship where maps are like lockers in particular and keys. There's a specific
data element that we keep track of in order to access another data type. So that's the
scope really of how maps are seen visually. I'm just giving a metaphor on that. So we
will use the maps to indicate what has been visited and the frequency of each
character. And the steps for us to be able to be able to create the sliding window
technique for this problems are going to be as follows.

Step one, I will also put it up here. We need to find the substring of S that includes
all the letters from T. So the steps for us to be able to solve this particular
problem are going to be as follows. Step one, we're going to initialize two pointers
left and right at the beginning of string S. Step two, we're going to initiate we're
going to initialize two frequency maps called frequency S and frequency T to store the
count of the characters in S and T. And we're going to create a variable counter to
keep track of the number of characters from t that are missing in the current window.

And we are going to create variables min length and start index to track the minimum
window length and the starting index of the minimum window substring. Then add the
frequency of characters from t to frequency t.
}

\expandafter\def\csname sourceexcerpt@48\endcsname{%
So pretty straightforward. We have an array of nums and we have an output of two. The
reason being is that with our following nums n= 3 since there are three numbers. So
all the numbers in the range 0 to three between the lowest point 0 and all the way to
the highest point we are accounting for all the numbers and we have to figure out
exactly which number is missing from our sequence between zero all the way to n and
we're going to be using bitwise manipulation for this technique bit manipulation That
is a freaking spell.

Anyway, so just to decide between our input and output, we are given an array of nums.
We need to return the missing number in the sequence that is between 0 and n. All
right. Right. Okay. So, the following steps for us to be able to accomplish this bit
manipulation is that we are going to use the sore operation used for exclusivity. And
the steps for us to accomplish this are going to be the following. Step one, we will
initialize a missing integer for the missing number. We will return.

Step two, we will calculate the or of all numbers from 0 to n and we will scan through
the array while doing so. We will sort the current number with the variable. And then
for the third step inside of our for loop or outside of it, we're going to calculate
the zor of all numbers in the given array. And while scanning each num we will use Zor
to update where the missing number is since we found the range. earlier between 0 to n
finally we will return our missing number. All right, just to give a breakdown I'm
just going to move this all up here for the sake of clarity.

So to go over the steps, we got step one, step two, step three, step four. And to make
this work, this is what we're going to do. Int missing number is going to be equal to
zero. Then what we will do is that we're going to scan for a range.
}

\expandafter\def\csname sourceexcerpt@49\endcsname{%
Also known as Hamming weight. Okay, that's pretty straightforward. They also give us a
couple notes and mention about Java and then also more Java, but we're working with
C++. So, well, it's pretty straightforward what they want us to do with the data. they
want us to write a function that takes the binary representation of an unsigned
integer and there are symbols that we use pretty much for the technique that we'll be
doing called bit manipulation and allows us to get the binary representation of our
following u for our following input where it looks like we're given into 32.

So we are given an unsigned integer and we need to return the number of one bits it
has. So looking at this, we just need to scan through all of our integers that we're
given for any well all the binaries with a given integer and we return the number of
ones. So the steps for us to be able to go about this and I will also have a visual at
the same time to be able to understand more on what forms of bit manipulation to use
which would be the following. So what we are going to do four steps are this.

Step one we have a given unsigned integer we will create a count to keep track of the
ones. Then what we're going to do is that while scanning through our unsigned integer
if n at the specified index, which is represent by and equals 1, we're going to
increase our count and then we will write shift the unsigned integer by one bit to
check for the rest of the bits and finally we will return our count. So the steps for
us to be able to go about this. Step one, step two, and then step three. Let me just
clean this up focusing on these steps we have over here.

So step one, the way how it's going to work, we're going to first create a count. Int
count is equal to zero. This is how we first initialize.
}
